\documentclass[11pt]{article}

\usepackage[margin=1in]{geometry}
\usepackage{amsmath,amssymb,bm}
\usepackage{booktabs}
\usepackage{hyperref}

\newcommand{\imc}{\textsc{imc}}
\newcommand{\ddmc}{\textsc{ddmc}}
\newcommand{\rich}{\textsc{rich}}
\newcommand{\storm}{\textsc{storm}}

\title{Why discrete-diffusion Monte Carlo stalls\\
the Giron converging Marshak wave}
\author{\rich{} / \storm{} implementation note}
\date{31 August 2026}

\begin{document}
\maketitle

\begin{abstract}
Enabling discrete-diffusion Monte Carlo (\ddmc) on Giron Test~2 produces a sphere that remains at the temperature floor.
The failure is not a lack of packets or a slow front.
It is the \imc--\ddmc{} interface at $r=R$.
The bath is optically thin and stays in \imc{}.
The unheated material is optically thick by eight orders of magnitude and is classified as \ddmc{}.
The Densmore--Wollaeger admission probability for that interface is $P\sim 4\times 10^{-9}$.
Rejected packets bounce back into the bath at unchanged weight and deposit nothing.
Ordinary \imc{} at the same face would absorb a few percent of the incident energy.
\ddmc{} therefore admits about $10^{-7}$ of the energy that \imc{} would, which is why the interior never heats.
\end{abstract}

\section{What was run}

Giron Test~2 is a spherical Marshak wave in a nonhomogeneous medium~\cite{giron}.
The material occupies $r<R$ with $R=0.05\,\mathrm{cm}$.
A thermostat surrounds $r\ge R$ and is reset every step to the paper's time-dependent bath temperature $T_{\mathrm{bath}}(t)$.
Hydrodynamics is off.
The mesh for job~10145269 used 128 radial shells and 8000 angular generators (1\,024\,000 interior cells).
The time step was $\Delta t=10^{-11}\,\mathrm{s}$.
Random walk was on; \ddmc{} was on.
By cycle~1800 ($t=-67.47\,\mathrm{ns}$ paper time) the analytic front sits at $r=0.0442\,\mathrm{cm}$, while 24 of 1\,024\,000 interior cells sit above the temperature floor, all in the single outermost radial shell.

The opacity, density, and material energy of Test~2 are
\begin{align}
  \kappa(T,\rho)
  &= 1.5\times 10^{4}\,
     \Bigl(\frac{T}{T_0}\Bigr)^{-3}
     \rho^{1.4}
     \qquad[\mathrm{cm}^{-1}],
  \label{eq:opacity}\\
  \rho(r)
  &= r^{1/2},
  \label{eq:density}\\
  u_m
  &= 3\times 10^{13}\,
     \Bigl(\frac{T}{T_0}\Bigr)^{2}
     \rho^{0.4}
     \qquad[\mathrm{erg}\,\mathrm{cm}^{-3}],
  \label{eq:eos}
\end{align}
with $T_0=0.1\,\mathrm{keV}$ (one hectoelectronvolt) and $\rho$ in $\mathrm{g}\,\mathrm{cm}^{-3}$.
There is no physical scattering: $\sigma_s=0$ and $\sigma_T=\sigma_a=\kappa$.
The temperature floor is $10^{-3}T_0$.
At $r=R$, $\rho=R^{1/2}=0.2236\,\mathrm{g}\,\mathrm{cm}^{-3}$.

\section{Two cells, two transport regimes}

A cell is treated as \ddmc{} when its mean-chord optical depth exceeds a threshold~\cite{densmore2012},
\begin{equation}
  \sigma_T\,\ell \ge \tau_{\mathrm{DDMC}},
  \qquad
  \ell = \frac{4V}{A},
  \qquad
  \tau_{\mathrm{DDMC}}=15
  \quad\text{(default)}.
  \label{eq:gate}
\end{equation}
The outermost interior cells have width $h=3.906\times 10^{-4}\,\mathrm{cm}$.
Taking $\ell\simeq 2h/3$ for a compact cell, Table~\ref{tab:two} follows from~\eqref{eq:opacity}.

\begin{table}[ht]
\centering
\caption{Transport quantities at $r=R$ for Giron Test~2 at cycle~1800 of job~10145269.
The bath temperature is $0.865\,T_0$.
The material is still at the floor $10^{-3}T_0$.}
\label{tab:two}
\begin{tabular}{@{}lcccc@{}}
\toprule
Region & $T/T_0$ & $\sigma_T$ [cm$^{-1}$] & $\lambda=1/\sigma_T$ [cm] & $\sigma_T\ell$ \\
\midrule
Bath ($r\ge R$) & $0.865$ & $2.85\times 10^{3}$ & $3.51\times 10^{-4}$ & $\sim 0.7$ \\
Material ($r<R$) & $10^{-3}$ & $1.84\times 10^{12}$ & $5.43\times 10^{-13}$ & $\sim 5\times 10^{8}$ \\
\bottomrule
\end{tabular}
\end{table}

The bath is below the gate~\eqref{eq:gate} and remains \imc{}.
Every unheated interior cell is above the gate by eight orders of magnitude and is classified as \ddmc{}.
The sphere surface $r=R$ is therefore an \imc$\to$\ddmc{} interface, not a \ddmc$\leftrightarrow$\ddmc{} face.

\section{The admission probability}

An \imc{} packet that hits a \ddmc{} face is not streamed into that cell.
\storm{} applies the static Densmore--Wollaeger asymptotic admission law~\cite{densmore2012,wollaeger2013}, implemented as \texttt{StaticAdmissionProbability} in \texttt{DDMCWollaegerInterface.hpp}:
\begin{equation}
  P_{\mathrm{admit}}(\mu)
  =
  \mathrm{clamp}\!\left(
    \frac{2\bigl(1+\tfrac{3}{2}\mu\bigr)}
         {3\bigl(\sigma_T d + z_0\bigr)}
  ,\, 0,\, 1\right).
  \label{eq:admit}
\end{equation}
Here $\mu=\bm{\Omega}\cdot\bm{n}$ is the cosine of the angle of incidence, $d$ is the target-cell centre-to-face distance, and $z_0=0.7104$ is the Milne extrapolation length in mean-free-path units.
If a uniform random number exceeds $P_{\mathrm{admit}}$, the packet undergoes a diffuse-albedo bounce in the \imc{} cell.
Its weight is unchanged.
No energy is tallied in the \ddmc{} cell.
The packet never becomes a \ddmc{} resident, so random walk in the material never runs.

The centre-to-face distance of the outermost interior shell is $d=h/2=1.953\times 10^{-4}\,\mathrm{cm}$.
In the cold material,
\begin{equation}
  \sigma_T d = (1.84\times 10^{12})(1.95\times 10^{-4}) = 3.59\times 10^{8}.
\end{equation}
The extrapolation length is then negligible, and~\eqref{eq:admit} collapses to
\begin{equation}
  P_{\mathrm{admit}}(\mu)
  =
  \frac{2+3\mu}{3\sigma_T d}.
  \label{eq:admit-thick}
\end{equation}
For normal incidence, $P_{\mathrm{admit}}(1)=4.64\times 10^{-9}$.
For isotropic incident intensity the flux-weighted mean is
\begin{equation}
  \langle P_{\mathrm{admit}}\rangle
  =
  \int_{0}^{1} P_{\mathrm{admit}}(\mu)\, 2\mu\,\mathrm{d}\mu
  =
  \frac{4}{3\sigma_T d}
  =
  3.71\times 10^{-9}.
  \label{eq:admit-mean}
\end{equation}
Three to four packets in a billion that strike $r=R$ from the bath are admitted.
The rest return to the bath.

This scaling is required by the diffusion boundary condition that~\eqref{eq:admit} encodes.
The same extrapolation length appears in the \ddmc$\to$\imc{} leak rate,
\begin{equation}
  \lambda_{\partial}
  =
  \frac{cA}{3V(\sigma_T d + z_0)}
  \propto
  \frac{1}{\sigma_T d}.
  \label{eq:leak}
\end{equation}
A cell with $\sigma_T d\sim 10^{8}$ both refuses incoming \imc{} packets and, if a packet were already inside, would leak them at a vanishing rate.

\section{What ordinary \imc{} would do at the same face}

In gray \imc{} the Fleck factor is
\begin{equation}
  f
  =
  \frac{1}{1 + 4aT^{3}\sigma_P c\Delta t / c_V},
  \label{eq:fleck}
\end{equation}
with $c_V=\partial u_m/\partial T$ from~\eqref{eq:eos}.
At the floor and $\Delta t=10^{-11}\,\mathrm{s}$, $f=1.09\times 10^{-3}$.
There is no physical scatter, so the only albedo is \imc{} effective scattering with single-scatter albedo $1-f$.
A diffusion estimate of the absorbed fraction for a semi-infinite medium is then
\begin{equation}
  \alpha_{\mathrm{IMC}}
  \simeq
  2\sqrt{\frac{f}{3}}
  =
  3.8\times 10^{-2}.
  \label{eq:imc-abs}
\end{equation}
About four percent of the energy that hits the cold material stays there.
The ratio of the two treatments at the same face is
\begin{equation}
  \frac{\langle P_{\mathrm{admit}}\rangle}{\alpha_{\mathrm{IMC}}}
  \simeq
  \frac{3.71\times 10^{-9}}{3.8\times 10^{-2}}
  \simeq
  1\times 10^{-7}.
  \label{eq:ratio}
\end{equation}
\ddmc{} is not a cheaper approximation of the same absorbed flux.
It suppresses the absorbed flux by seven orders of magnitude.

The origin of the mismatch is physical.
Equation~\eqref{eq:admit} is an asymptotic diffusion boundary condition derived for a scattering-dominated, nearly isotropic interior~\cite{densmore2012,wollaeger2013}.
The cold Marshak material is the opposite: $\sigma_s=0$, $\sigma_a$ is huge, and $f\ll 1$.
A photon that enters is absorbed in $\lambda=5\times 10^{-13}\,\mathrm{cm}$, a distance $10^{9}$ times smaller than the cell.
The correct transmission is set by the Fleck-effective albedo~\eqref{eq:imc-abs}, not by $1/(\sigma_T d)$.

\section{What the run actually did}

Thirty-seven million packets are generated in the bath each step.
Over 1800 steps that is enough for a handful of admissions even at $P\sim 4\times 10^{-9}$.
Each admitted packet carries full bath weight into one cell and can raise that cell from the floor to $\sim 0.87\,T_0$ in a single event.
That is the observed pattern: 0, 4, 6, 8, 10, 11, 17, 22, then 24 cells above the floor at cycles 200 through 1800, all at $r=0.04980\,\mathrm{cm}$, with a bimodal temperature (floor almost everywhere, bath temperature in those 24 cells).

The process is self-locking.
Once a cell reaches $0.87\,T_0$ its opacity drops to $\sigma_T\simeq 2.85\times 10^{3}\,\mathrm{cm}^{-1}$, $\sigma_T\ell\sim 1$, and it leaves \ddmc{}.
Its inward neighbours remain at the floor with $P_{\mathrm{admit}}\sim 10^{-9}$.
Nothing propagates inward.
The deepest heated radius is therefore frozen from cycle~400 through cycle~1800.

Raising the floor does not repair the law.
At $T=0.1\,T_0$, $\sigma_T d\sim 360$ and $\langle P_{\mathrm{admit}}\rangle\sim 4\times 10^{-3}$, still an order of magnitude below $\alpha_{\mathrm{IMC}}$.
Only when the target cell is already hot enough to leave \ddmc{} does admission become order unity, which is exactly when \ddmc{} is no longer applied.

\section{A second, independent stall inside the sphere}

Suppose a packet is admitted, or that two neighbouring material cells are both \ddmc{}.
The internal leak rate between \ddmc{} cells $i$ and $j$ is~\cite{densmore2012}
\begin{equation}
  R_{ij}
  =
  \frac{d_{iF}}{D_i}+\frac{d_{jF}}{D_j},
  \qquad
  D=\frac{c}{3\sigma_T},
  \qquad
  \lambda_{i\to j}
  =
  \frac{A_F}{V_i R_{ij}}.
  \label{eq:internal}
\end{equation}
If cell $i$ is heated and cell $j$ is at the floor, $D_j/D_i=\sigma_i/\sigma_j\sim 10^{-9}$.
The resistance is then dominated by the cold cell, $R_{ij}\simeq d_{jF}/D_j$, and
\begin{equation}
  \lambda_{i\to j}
  \simeq
  \frac{A_F D_j}{V_i d_{jF}}
  \propto
  \frac{1}{\sigma_j}.
\end{equation}
Leakage into the unheated neighbour is diffusion-limited at the cold opacity.
This is the stall already noted in the driver comment in \texttt{test.cpp}:
the leak rate into an unheated cell whose optical depth is many orders of magnitude larger than the heated cell behind the front cannot carry the incident bath flux.
Plain \imc{} with random walk does, because an \imc{} packet that has entered a cold cell is absorbed locally on the Fleck time scale rather than waiting for a diffusion leak.

Either mechanism is fatal.
The interface law~\eqref{eq:admit} prevents energy from crossing $r=R$.
The internal law~\eqref{eq:internal} would prevent a front from advancing even if energy did cross.

\section{Why the existing \ddmc{} parameters cannot fix this}

The configurable fields are $\tau_{\mathrm{DDMC}}$, a particle-depth gate, an external-source face gate, and moving-interface corrections.
None of them appears in~\eqref{eq:admit}.
The extrapolation length $z_0=0.7104$ is a compile-time constant.
The product $\sigma_T d$ is set by the opacity law~\eqref{eq:opacity} and the mesh.

Lowering $\tau_{\mathrm{DDMC}}$ makes more cells \ddmc{} and cannot increase $P_{\mathrm{admit}}$.
Raising $\tau_{\mathrm{DDMC}}$ above $\sigma_T\ell\sim 5\times 10^{8}$ in the cold material turns \ddmc{} off in the region that matters, which is the same as \texttt{--ddmc false}.
The moving-interface factor $G_U(\mu)$ is irrelevant: the mesh is static.

\section{Conclusion}

On Giron Test~2 the \imc--\ddmc{} interface at $r=R$ has admission probability $3.7\times 10^{-9}$ against an \imc{} absorbed fraction of $3.8\times 10^{-2}$.
That single number explains the stalled run: the bath is hot and well resolved, packets are plentiful, and the interior remains at the floor because the interface reflects them.
The asymptotic admission law is the wrong boundary condition for a purely absorbing, Fleck-scattered Marshak toe.
Until a different interface treatment is derived for absorption-dominated media, this benchmark should be run with \imc{} and random walk only.

\begin{thebibliography}{9}
\bibitem{giron}
J.~F. Giron et~al.,
\textit{Converging Marshak waves in nonhomogeneous media}.

\bibitem{densmore2012}
J.~D. Densmore, K.~G. Thompson, and T.~J. Urbatsch,
``A hybrid transport-diffusion Monte Carlo method for frequency-dependent radiative-transfer simulations,''
\textit{J.\ Comput.\ Phys.} \textbf{231}, 6924--6934 (2012).

\bibitem{wollaeger2013}
R.~T. Wollaeger, D.~R. van~Rossum, C.~Graziani, S.~M. Couch, G.~C. Jordan~IV, D.~Q. Lamb, and G.~A. Moses,
``Radiation transport for explosive outflows: a multigroup hybrid Monte Carlo method,''
\textit{Astrophys.\ J.\ Suppl.\ Ser.} \textbf{209}, 36 (2013).
\end{thebibliography}

\end{document}
